% Generated from validated Article Draft Result. Do not hand-edit; revise the structured draft instead.
\section{生成が開発環境へ入る――CodexからCopilot previewへ}
\label{pkg:code-developer-interface}
\noindent\textbf{CodexとGitHub Copilotを同じ『コード生成AI』として一つの出来事にまとめると、2021年の変化は見えにくい。Codexはcode-generation modelの研究・評価とlifecycleを示し、Copilot technical previewはその能力がdeveloper workflowのinterfaceへ組み込まれる別の境界を示した。model capabilityとproduct interfaceの間に、生成AIの新しいdistribution layerが生まれた。}\autocite{src-089fe2c6fbb1,src-47d0c87b08f0,src-5a992319f41e,src-890a73c0b288}

Codexの一次資料は、large language modelをcodeへ適用するtechnical evaluationと、その後の公開surfaceを含むlifecycleを2021年の記録として残している。ここで重要なのは、単に『コードも書けるようになった』という能力追加ではない。自然言語生成の延長にcode-specialized generationが明確な研究対象として立ち上がり、model behaviorをsoftware developmentという具体的な作業領域で評価する枠組みが前景化したことである。\autocite{src-47d0c87b08f0,src-5a992319f41e,src-890a73c0b288}


GitHub Copilot technical previewは、研究modelとは別のproduct boundaryである。developerが日常的に使う開発環境の中へ生成機能を持ち込み、modelへ直接promptすること自体を主役にせず、作業interfaceの一部として能力を提示する。この違いは大きい。研究成果が存在することと、利用者のworkflowに埋め込まれたproduct experienceとして提供されることの間には、interface designとdistributionの層がある。\autocite{src-089fe2c6fbb1}


したがってCodexとCopilotの関係は、同じartifactの別名ではなく、model capabilityとproduct interfaceという異なるレイヤの接続として読むべきである。model側では何を生成できるか、どのように評価されるかが中心になる。一方product側では、利用者がどこでその能力に触れ、既存の作業手順の中でどう呼び出すかが中心になる。2021年には、生成AIの価値がmodel単体のbenchmarkからworkflow integrationへ広がる兆候がはっきり現れた。\autocite{src-089fe2c6fbb1,src-47d0c87b08f0,src-5a992319f41e,src-890a73c0b288}


この章ではpaper、research model、product previewを同じchronology stateとして扱わない。Codexのtechnical evaluation、Codexとしてのmodel/service lifecycle、Copilot technical previewは、互いに関連していても別の公開境界である。後年のCopilot製品の成熟を2021年previewへ逆投影せず、当時確認できるのはcode generationが研究成果からdeveloper-facing interfaceへ接続し始めた段階までとする。\autocite{src-089fe2c6fbb1,src-47d0c87b08f0,src-5a992319f41e,src-890a73c0b288}


年次全体では、このdeveloper interfaceはaccess/distributionの章とも響き合う。API、open weights、product previewはいずれもmodelの外側にあるが、誰がどの場面で能力へアクセスできるかを決める。CodexからCopilotへの接続は、その中でも『開発者が既にいる場所へ生成AIが入る』という新しいsurfaceを示した。2022年以降のtool useやagentを先取りしたと断定する必要はなく、まずworkflow-embedded generationの初期境界として評価するのが妥当である。\autocite{src-089fe2c6fbb1,src-47d0c87b08f0,src-5a992319f41e,src-890a73c0b288}


\begin{claimboundary}
Codexのtechnical evaluationやCopilotのproduct capabilityに関する評価は、それぞれauthor/project側の主張として扱う。一次資料の確認は、性能・生産性・安全性を独立再現したことを意味しない。\autocite{src-089fe2c6fbb1,src-47d0c87b08f0,src-5a992319f41e,src-890a73c0b288}
\end{claimboundary}
